% =====================================================================
% Appendix: proofs for Section 6 (Coherence: what contradictions fix)
% Sources: research/theory/T2-coherence-as-negative-data.md (as repaired; see its
%          "## Verification log"), research/verification/T2-...-verification.md
%          (referee countermodels and checks), research/lit/L3-...md (Thm D, Thm D'),
%          research/lit/L5-logic-facts.md (Prop 2.5).
% =====================================================================
\providecommand{\Steps}{\mathrm{Steps}}
\providecommand{\Taut}{\mathrm{Taut}}
\providecommand{\Cai}{C_{\mathrm{ai}}}
\providecommand{\CFm}{C_{\Fm}}
\providecommand{\Cadm}{C_{\mathrm{adm}}}
\providecommand{\BV}{\mathrm{BV}}
\providecommand{\TTs}{\mathrm{TT}}
\providecommand{\RobQ}{\mathsf{Q}}
\providecommand{\valtop}{v_{\top}}
\providecommand{\valtaut}{v_{\Taut}}
\providecommand{\hstar}{h^{*}}
\providecommand{\dstar}{\der^{*}}

\section{Proofs for Section~\ref{sec:coherence}}\label{app:coherence}

\Cref{lem:coherence:bag}, \cref{prop:coherence:caution}(a),(c), \cref{thm:coherence:halving}, \cref{thm:coherence:robust}, \cref{thm:coherence:silent}, \cref{prop:coherence:chains}, \cref{cor:coherence:limitpoints}, \cref{lem:coherence:delta0}, \cref{thm:coherence:turing}, \cref{thm:coherence:trilemma}, \cref{thm:coherence:carnap} (given \cref{lem:coherence:duality}), \cref{thm:coherence:rank}, \cref{prop:coherence:tonk}, \cref{thm:coherence:elimination}, \cref{prop:coherence:quus} (immediate), \cref{cor:coherence:fallacies} and \cref{lem:coherence:projection} are proved in the main text, and \cref{thm:coherence:rosser} is cited. This appendix proves the remaining results. The proofs follow T2 in the form repaired after adversarial verification; where a statement or proof was changed by the repair, we say so. Notation is as in \cref{sec:coherence}. For formulas, $\equiv_2$ is classical equivalence.

% ---------------------------------------------------------------------
\subsection{Proof of \texorpdfstring{\cref{prop:coherence:oligarchy}}{the oligarchy proposition}}\label{app:coherence:oligarchy}

($\Leftarrow$) Let $\hat R\subseteq\bigcap S$ with $w(S)\ge\frac12w(\VS)$. A detection has $\Steps(\pi)=P\subseteq\hat R$, so every $R\in S$ contains $P$ and is deleted; the deleted weight is at least $w(S)$.

($\Rightarrow$) The set of steps is countable, because $\Fm$ is. If $\hat R=\emptyset$, take $S=\VS$. Otherwise enumerate $\hat R=\{s_1,s_2,\dots\}$ (finite or infinite), and let $P_k=\{s_1,\dots,s_k\}$ and $S_k=\{R\in\VS:P_k\subseteq R\}$. If $P_k\not\subseteq\bigcap\VS$, then by the worst-case assumption $P_k$ is a possible detection, and it deletes exactly $S_k$; so $w(S_k)\ge\frac12w(\VS)$. If $P_k\subseteq\bigcap\VS$, then $S_k=\VS$. In both cases $w(S_k)\ge\frac12w(\VS)$. The sets $S_k$ decrease, and $\bigcap_kS_k=\{R\in\VS:\hat R\subseteq R\}=:S$. Since $w$ is a finite measure on the countable set $\VS$, continuity from above gives $w(S)=\lim_kw(S_k)\ge\frac12w(\VS)>0$, and $\hat R\subseteq\bigcap S$ by definition. (For finite $\VS$ one may instead take one step $s_R\in\hat R\setminus R$ for each $R$ with $\hat R\not\subseteq R$; this is the version formalized in Lean.) \qed

\emph{Repair note.} The first version assumed $\VS$ finite and allowed every finite $P\subseteq\hat R$ as a detection, including $P\subseteq\bigcap\VS$, which would delete the target and contradicts \cref{lem:coherence:bag}; the hypothesis is now restricted to legal detections and $\VS$ may be countable.

% ---------------------------------------------------------------------
\subsection{Proofs of \texorpdfstring{\cref{prop:coherence:caution}(b) and \cref{prop:coherence:tradeoff}}{the caution and trade-off propositions}}\label{app:coherence:tradeoff}

\paragraph{\Cref{prop:coherence:caution}(b).} Let $s_c:=\bigvee_i(a_i\to\ell_i^{1-c_i})$ for $c\in\{0,1\}^k$. If $b_i\neq c_i$ for some $i$, then $\ell_i^{1-c_i}=\ell_i^{b_i}$, so the disjunct $a_i\to\ell_i^{b_i}$ of $s_c$ is an axiom of $h_b$ and $\step s_c\in h_b$. If $b=c$, the valuation with every $a_i$ true and $p_i=c_i$ (so $\ell_i^{c_i}$ true and $\ell_i^{1-c_i}$ false) satisfies the axioms of $h_c$ and falsifies $s_c$. So $\step s_c\in h_b$ iff $b\neq c$. Every $h_b$ is consistent, so $\Ac=\{\emptyset\}$ is coherent for all of them. A certified-sound learner accepts only steps in $\bigcap\VS_t$, so it rejects $\step s_c$ while $h_c\in\VS_t$. For target $h_{b^*}$, present $\step s_c$ for each $c\neq b^*$ in any order: each is valid in the target, each deletes exactly $h_c$, and at the moment it is presented $h_c$ has not yet been deleted, so each is an incompleteness error. That makes $2^k-1$ errors, although $\log_2|\Hc|=k$. \qed

\paragraph{\Cref{prop:coherence:tradeoff}.} \emph{(a)} As just shown, $\step s_c\in h_b$ iff $b\neq c$. Fix a learner that in every round announces $\hat R_t\subseteq R$ for some $R\in\VS_t$; we may assume it is deterministic, and the environment sees $\hat R_t$ before it moves. While $|\VS_t|\ge2$, the adversary picks $h_c\in\VS_t$ with $\hat R_t\subseteq h_c$ and presents $\step s_c$. Since ${\step s_c}\notin h_c\supseteq\hat R_t$, this is an incompleteness error (for any target other than $h_c$). Since $\step s_c\in h_b$ for every $b\neq c$, it deletes exactly $h_c$. After $2^k-1$ rounds a single hypothesis $h_{b^*}$ remains; take it as the target. Every presented datum $\step s_c$ has $c\neq b^*$, hence lies in $h_{b^*}$, so the run is a prefix of a text for $h_{b^*}$; extend it arbitrarily. The target must be chosen at the end: if $\VS_t=\{h_c,h_{b^*}\}$ and the learner's coalition is $\{h_{b^*}\}$, a target fixed in advance as $h_{b^*}$ would not be charged the last error.

\emph{The union learner.} If $\Rstar\in\VS_t$ then $\Rstar\subseteq\bigcup\VS_t=\hat R_t$, so no incompleteness error occurs; and $\Rstar$ is never deleted, by \cref{lem:coherence:bag}. Let $h_\neg=\{\Gamma\step\varphi:\ \Gamma\cup\{\neg a_1,\dots,\neg a_k\}\models\varphi\}$. Since $\neg a_i\models a_i\to\ell$, every $h_b\subseteq h_\neg$, so $\hat R_t\subseteq h_\neg$ and $\Cl_{\hat R_t}(\emptyset)\subseteq\Cn_2(\{\neg a_1,\dots,\neg a_k\})$, which does not contain $\bot$. So no argument from $\emptyset$ to $\bot$ uses only accepted steps: there are no detections. For target $h_{b^*}$, as long as $\VS_t$ contains some $h_b\neq h_{b^*}$ the union learner accepts $\step a_i\to\ell_i^{b_i}$ for an $i$ with $b_i\neq b^*_i$, which is not in $h_{b^*}$ (make $a_i$ and $\ell_i^{b^*_i}$ true); so it is unsound, and it is not oligarchic, so \cref{thm:coherence:halving} does not apply to it.

\emph{(b)} Let $\Hc=\{h_1,h_2,h_3\}$ be the class of \cref{thm:search:doctrinal}, each with weight $\frac13$. Any step $s\in\bigcup\VS\setminus\hat R$ can be presented as a positive datum for a target containing it, and then it is an incompleteness error that deletes $\{h_i:s\notin h_i\}$. This weighs at least $\frac12$ only if at least two $h_i$ omit $s$, i.e.\ only if $s$ lies in at most one $h_i$. So if every possible incompleteness error deletes at least half the weight, $\hat R\supseteq M:=\{s:\ s\in h_i\text{ for at least two }i\}$. But the argument $\pi$ of \cref{thm:search:doctrinal} has $\Steps(\pi)\subseteq M\subseteq\hat R$ and $\Steps(\pi)\not\subseteq h_i$ for each $i$. So it is a legal detection for every target, and it deletes nothing. Equivalently, by \cref{prop:coherence:oligarchy}, detection-halving needs $\hat R\subseteq h_i\cap h_j$ for some $i\neq j$, and $M\not\subseteq h_i\cap h_j$ for every pair: $h_1\cap h_2$ lacks $\step q$, while $h_1\cap h_3$ and $h_2\cap h_3$ lack $\step p$. \qed

\emph{Repair note.} Both parts were added after verification. The first version concluded from \cref{thm:coherence:halving} that ``coherence bounds the price of boldness while nothing bounds the price of caution''; part (a) shows that on \cref{prop:coherence:caution}(b)'s own class every oligarchic learner pays the price of caution. Membership facts were checked by truth tables for $k\le4$ (T2 verification log).

% ---------------------------------------------------------------------
\subsection{Proof of \texorpdfstring{\cref{prop:coherence:informant}}{the informant proposition}}\label{app:coherence:informant}

\emph{(a)} For the target, $\Gamma\nder^*\varphi$ iff $\Gamma\cup\{\neg\varphi\}$ is target-coherent, by classical reductio. So the designations eventually enumerate (codes for) every finite-premise sequent not in $\dstar$, and the text enumerates $\dstar$: together they form an informant. Let $(h_i)$ enumerate $\Hc$, and let the learner output the least $i$ such that $h_i$ contains every sequent seen in the text and is coherent on every designated context seen. Let $i^*$ be the least index of the target and $j<i^*$, so $h_j\neq{\dstar}$. If ${\dstar}\not\subseteq h_j$, the text eventually shows a sequent of ${\dstar}\setminus h_j$, and $j$ is rejected from then on. Otherwise pick $\Gamma\step\varphi\in h_j\setminus{\dstar}$. Then $\Gamma\cup\{\neg\varphi\}$ is target-coherent, hence eventually designated, and it is $h_j$-incoherent by reductio in $h_j$; so $j$ is rejected from then on. The target's index $i^*$ is never rejected. As there are finitely many $j<i^*$, the output is eventually $i^*$ forever. If membership in the $h_i$ is uniformly decidable, both tests are decidable (coherence of $h_i$ on $A$ is non-membership of $A\step\bot$), so the learner is computable.

\emph{(b)} $\{p_j\}$ is $h_n$-coherent iff $\{\neg p_i:i<n\}\cup\{p_j\}$ is consistent iff $j\ge n$, and it is $h_\omega$-incoherent for every $j$. The chain is strictly increasing with union $h_\omega$, and all members are $\emptyset$-coherent. Consider the computable learner that outputs $h_j$ for the least $j$ such that $\{p_j\}$ has been designated, and $h_\omega$ if there is none (the enumeration learner of (a) also works). On target $h_n$, exactly the singletons $\{p_j\}$ with $j\ge n$ can be designated, and $\{p_n\}$ eventually is; from then on the least designated $j$ is $n$, and the learner outputs $h_n$. On target $h_\omega$ no singleton $\{p_j\}$ is ever designated, and the learner outputs $h_\omega$ throughout. \qed

\emph{Repair note.} This proposition was added after verification; the referee's counterexample (b) showed that the first version of \cref{cor:coherence:limitpoints}, stated without ``$\Ac$ fixed'', was false.

% ---------------------------------------------------------------------
\subsection{Proofs for \texorpdfstring{\cref{sec:coherence:post}}{Section 6.2}: Post-completeness}\label{app:coherence:post}

\paragraph{\Cref{thm:coherence:post}, pure $\{\to\}$.} Let $\varphi\in C(X)\setminus\Cn_2(X)$ and $v$ a Boolean valuation with $v[X]=1$, $v(\varphi)=0$. Put $\sigma_v(p)=p_0\to p_0$ if $v(p)=1$ and $\sigma_v(p)=p_0$ otherwise. By induction on $\psi$, $\sigma_v\psi\equiv_2p_0\to p_0$ if $v(\psi)=1$ and $\sigma_v\psi\equiv_2p_0$ if $v(\psi)=0$. For atoms this is the definition. For $\psi=\alpha\to\beta$: if $v(\alpha)=0$, then $\sigma_v\psi\equiv_2p_0\to(\cdot)$, and both $p_0\to(p_0\to p_0)$ and $p_0\to p_0$ are tautologies; if $v(\alpha)=v(\beta)=1$, then $\sigma_v\psi\equiv_2\top\to\top\equiv_2\top$; if $v(\alpha)=1$, $v(\beta)=0$, then $\sigma_v\psi\equiv_2\top\to p_0\equiv_2p_0$. Hence $\sigma_vX\subseteq\Taut\subseteq C(\emptyset)$ and, by structurality, $\sigma_v\varphi\in C(\sigma_vX)\subseteq C(\emptyset)$, with $\sigma_v\varphi\equiv_2p_0$. So $p_0\in\Cn_2(\{\sigma_v\varphi\})\subseteq C(C(\emptyset))=C(\emptyset)$. Substituting any $\psi$ for $p_0$ and using $\sigma C(\emptyset)\subseteq C(\emptyset)$ gives $C(\emptyset)=\Fm$. \qed

\paragraph{\Cref{prop:coherence:almost}.} $\Cai$ is a closure operator, and it is structural because $\sigma X$ is non-empty iff $X$ is. In the $\{\wedge,\vee\}$ language $\Cn_2(\emptyset)=\emptyset$, since the all-false valuation falsifies every formula; so $\Cai\supseteq\Cn_2$. Now let $C\supsetneq\Cn_2$ be structural.
\emph{Case $C(\emptyset)\neq\emptyset$.} Pick $\varphi\in C(\emptyset)$ and substitute $p_0$ for every atom. Every one-variable lattice term is $\equiv_2p_0$, so $p_0\in\Cn_2(\{\sigma\varphi\})\subseteq C(C(\emptyset))=C(\emptyset)$, and by structurality $C(\emptyset)=\Fm$: $C=\CFm$.
\emph{Case $C(\emptyset)=\emptyset$.} Take $\varphi\in C(X)\setminus\Cn_2(X)$ (so $X\neq\emptyset$) and a Boolean $v$ with $v[X]=1$, $v(\varphi)=0$. Let $\sigma$ send the atoms true under $v$ to $q$ and the false ones to $r$ (distinct atoms). Each $\sigma\psi$ is a lattice term in $q,r$, and its value at $(q,r)=(1,0)$ is $v(\psi)$. Up to $\equiv_2$, the lattice terms in $q,r$ are $q$, $r$, $q\wedge r$ and $q\vee r$. So $\sigma\psi\in\{q,\,q\vee r\}$ up to $\equiv_2$ for $\psi\in X$, and $\sigma\varphi\in\{r,\,q\wedge r\}$. Hence $\sigma X\subseteq\Cn_2(\{q\})$ and $r\in\Cn_2(\{\sigma\varphi\})$, so
\[
r\in C(\{\sigma\varphi\})\subseteq C(C(\sigma X))=C(\sigma X)\subseteq C(\Cn_2(\{q\}))\subseteq C(C(\{q\}))=C(\{q\}).
\]
Substituting $\chi$ for $q$ and $\psi$ for $r$ gives $\psi\in C(\{\chi\})$ for all $\chi,\psi$, so $C(Y)=\Fm$ for every non-empty $Y$: $C=\Cai$. \qed

\paragraph{\Cref{thm:coherence:cpc}.} \emph{(a)} Suppose $\langle D\rangle=\Cn_2$. Every $h\in\mathcal S(D,A_0)$ is structural and contains $D$, hence contains $\langle D\rangle=\Cn_2$; by \cref{thm:coherence:post}, $h\in\{\Cn_2,\CFm\}$, and $A_0$ excludes $\CFm$. Conversely, suppose $\langle D\rangle\subsetneq\Cn_2$ (note $\langle D\rangle\subseteq\Cn_2$ because $\Cn_2$ is structural and contains $D$). Then $\langle D\rangle$ is structural, contains $D$, and is $A_0$-coherent because $\Cn_2$ is; so $\langle D\rangle$ and $\Cn_2$ are distinct members of $\mathcal S(D,A_0)$.

\emph{Explicit finite bases.} For a language containing $\neg$ and $\to$ (possibly also $\wedge,\vee$), with distinct atoms $p,q,r$, take the axioms $\step p\to(q\to p)$, $\step(p\to(q\to r))\to((p\to q)\to(p\to r))$, $\step(\neg p\to\neg q)\to(q\to p)$ and modus ponens $p,\ p\to q\step q$; if $\wedge,\vee$ are present, add their clauses \emph{as $\to$-axioms}: $\step p\wedge q\to p$, $\step p\wedge q\to q$, $\step p\to(q\to p\wedge q)$, $\step p\to p\vee q$, $\step q\to p\vee q$, $\step(p\to r)\to((q\to r)\to(p\vee q\to r))$. For pure $\{\to\}$ take the first two axioms, Peirce's law $\step((p\to q)\to p)\to p$ and modus ponens. For languages without $\to$, such as $\{\neg,\wedge\}$ and $\{\neg,\vee\}$, finite bases exist because every two-element matrix logic is finitely based \citep[as we recall the result; we have not re-checked its exact statement]{rautenberg1981two}.

The least structural consequence relation generated by a set of sequents is given by derivations from substitution instances of those sequents. With the axioms and modus ponens above, these are derivations in a Hilbert system with modus ponens as its only rule: \L ukasiewicz's axioms for $\{\neg,\to\}$, extended by the usual axioms for $\wedge$ and $\vee$, and for $\{\to\}$ the axioms $K$, $S$ and Peirce's law, a variant of the Tarski--Bernays axiomatization ($K$ and $S$ give intuitionistic implication, and Peirce's law makes it classical). These systems are complete for tautologies \citep{kleene1952introduction,mendelson1964introduction}; for the $\wedge,\vee$ axioms Kalm\'ar's lemma goes through because the system proves $p,q\vdash p\wedge q$, $\neg p\vdash\neg(p\wedge q)$, $\neg q\vdash\neg(p\wedge q)$, $p\vdash p\vee q$, $q\vdash p\vee q$ and $\neg p,\neg q\vdash\neg(p\vee q)$. Since modus ponens is the only rule and the first two axioms are present, the deduction theorem holds. So if $X\models\varphi$, compactness gives a finite $\{x_1,\dots,x_n\}\subseteq X$ with $\models x_1\to(\cdots\to(x_n\to\varphi))$, which is provable, and $n$ applications of modus ponens derive $\varphi$ from $X$. Hence $\langle D\rangle=\Cn_2$.

\emph{(b)} Let $j<i^*$, so $h_j\neq\Cn_2$. If $h_j\not\supseteq\Cn_2$, some sequent of $\Cn_2\setminus h_j$ eventually appears in the text, and $j$ is ineligible from then on. If $h_j\supseteq\Cn_2$, then $h_j=\CFm$ by \cref{thm:coherence:post}, and $A_0\der_{h_j}\bot$ makes $j$ ineligible from the start. $i^*$ is always eligible. So the learner's output is non-decreasing, never exceeds $i^*$, and is eventually $i^*$; in particular it changes its mind at most $i^*$ times. Uniform decidability makes the eligibility tests computable; no property of the ordering was used.

\emph{(c)} $\Cn_2\oplus\{\step p\}$ with $p$ not in $A_0$: its consequences of $A_0$ are $\Cn_2(A_0\cup\{p\})$, which is consistent, and it contains every sequent of $\Cn_2$. $\CFm$ contains every sequent. $\Cai$ contains every sequent of $\Cn_2$ in the $\{\wedge,\vee\}$ fragment and is non-trivial on $\emptyset$. \qed

\paragraph{The sequent-rule basis fails (countermodel).} Take the language $\{\neg,\to,\wedge\}$ and let $D'$ consist of the three $\{\neg,\to\}$ axioms, modus ponens, and the sequent rules $p\wedge q\step p$, $p\wedge q\step q$, $p,q\step p\wedge q$. We show $\step p\to(q\to p\wedge q)\notin\langle D'\rangle$; this countermodel is due to the verifier of T2 and was re-checked by script. For every assignment $s$ of truth values to atoms define $v_s$ on all formulas, simultaneously for all $s$, by recursion on complexity: $v_s(p)=s(p)$; $v_s$ treats $\neg$ and $\to$ classically; and $v_s(\alpha\wedge\beta)=1$ iff $\alpha\in\Theta^*$ and $\beta\in\Theta^*$, where $\Theta^*=\{\varphi:\ v_s(\varphi)=1\text{ for all }s\}$. This is well defined, since membership of $\alpha,\beta$ in $\Theta^*$ depends only on values at lower complexity. Every substitution instance of a $\{\neg,\to\}$ tautology lies in $\Theta^*$, because each $v_s$ evaluates $\neg,\to$ classically and treats the substituted subformulas as atoms. $\Theta^*$ is closed under modus ponens (if $v_s(\alpha)=v_s(\alpha\to\beta)=1$ for all $s$, then $v_s(\beta)=1$ for all $s$), under $\wedge$-introduction (if $\alpha,\beta\in\Theta^*$ then $v_s(\alpha\wedge\beta)=1$ for all $s$) and under $\wedge$-elimination (if $\alpha\wedge\beta\in\Theta^*$ then $\alpha,\beta\in\Theta^*$ by definition). So $\langle D'\rangle(\emptyset)\subseteq\Theta^*$. But $p\notin\Theta^*$, so $v_s(p\wedge q)=0$ for every $s$; at $s(p)=s(q)=1$ this gives $v_s(p\to(q\to p\wedge q))=0$. Hence $p\to(q\to p\wedge q)\notin\Theta^*$, and $\langle D'\rangle\neq\Cn_2$. \qed

\paragraph{Renaming is not enough (L5 Prop 2.5).} Let $C(X)=\Cn_2(X)$ if $\Cn_2(X)$ contains no atom, and $C(X)=\Cn_2(X\cup\mathrm{At})$ otherwise, where $\mathrm{At}$ is the set of atoms. $C$ is extensive; it is monotone because $\Cn_2(X)\subseteq\Cn_2(Y)$ for $X\subseteq Y$; and it is idempotent, since in the first case $\Cn_2(\Cn_2(X))$ has no atom and in the second $C(X)$ already contains every atom. For a renaming $\rho$ (atoms to atoms) we have $\rho C(X)\subseteq C(\rho X)$: in the first case because $\Cn_2$ is structural, and in the second because $\Cn_2(\rho X)$ contains $\rho p$ for an atom $p\in\Cn_2(X)$, so $C(\rho X)=\Cn_2(\rho X\cup\mathrm{At})\supseteq\rho\Cn_2(X\cup\mathrm{At})$. $C$ contains $p\step q$, and $C(\emptyset)=\Taut$; so $C$ is $\emptyset$-coherent, has the classical theorems, and differs from $\Cn_2$. It is the least renaming-invariant closure operator above $\Cn_2$ containing $p\step q$: any such operator puts every atom into $C'(X)$ as soon as $\Cn_2(X)$ contains one. \qed

\paragraph{\Cref{prop:coherence:adm}(a).} \emph{Extensive and monotone:} clear. \emph{Idempotent:} if $\sigma X\subseteq\mathrm{Th}$, then $\sigma\Cadm(X)\subseteq\mathrm{Th}$ by definition, so every $\varphi\in\Cadm(\Cadm(X))$ has $\sigma\varphi\in\mathrm{Th}$. \emph{Structural:} if $\varphi\in\Cadm(X)$ and $\sigma(\tau X)\subseteq\mathrm{Th}$, then applying the definition to $\sigma\circ\tau$ gives $\sigma(\tau\varphi)\in\mathrm{Th}$; so $\tau\varphi\in\Cadm(\tau X)$. $\Cadm(\emptyset)=\mathrm{Th}$ because $\mathrm{Th}$ is substitution-closed. \emph{Maximal:} if $C$ is structural with $C(\emptyset)=\mathrm{Th}$, $\varphi\in C(X)$ and $\sigma X\subseteq\mathrm{Th}$, then $\sigma\varphi\in C(\sigma X)\subseteq C(C(\emptyset))=\mathrm{Th}$. For finite $X$, the defining condition is admissibility of $X/\varphi$. \qed

\paragraph{\Cref{prop:coherence:adm}(b),(c).} \emph{(b)} $\Cn_2$ is structural with theorems $\Taut$, so $\Cn_2\subseteq\Cadm$ by (a). If $\varphi\notin\Cn_2(X)$, the substitution $\sigma_v$ of the proof of \cref{thm:coherence:post} maps $X$ into $\Taut$ and $\varphi$ to a contradiction (to a formula $\equiv_2p_0$ in the pure $\{\to\}$ case), which is not in $\Taut$; so $\varphi\notin\Cadm(X)$. In the $\{\wedge,\vee\}$ fragment $\mathrm{Th}=\emptyset$, so the condition ``$\sigma X\subseteq\mathrm{Th}$'' never holds for non-empty $X$: $\Cadm(X)=\Fm$ for $X\neq\emptyset$ and $\Cadm(\emptyset)=\emptyset$, i.e.\ $\Cadm=\Cai$.

\emph{(c)} The admissibility and non-derivability of Harrop's rule are cited \citep{harrop1960concerning,rybakov1997admissibility,iemhoff2001admissible}. The two relations have the same theorems by (a). For the coherence profile, first note that $\IPC$ decides every variable-free formula according to its classical value: by induction, $\vdash_{\IPC}\neg\bot$; and if $\alpha,\beta$ are each provable or refutable according to their classical values, so are $\alpha\wedge\beta$, $\alpha\vee\beta$, $\alpha\to\beta$ and $\neg\alpha$. Now if $A$ is classically consistent, pick a Boolean $v\models A$ and substitute $\top:=\bot\to\bot$ or $\bot$ for each atom according to $v$. Then $\sigma_vA$ consists of classically true variable-free formulas, all $\IPC$-theorems, while $\sigma_v\bot=\bot$ is not one; so the rule $A/\bot$ is not admissible, i.e.\ $A\nder^{\IPC}_{\mathrm{adm}}\bot$. If $A$ is classically inconsistent, then $A\der_{\IPC}\bot$ by Glivenko's theorem (\cref{thm:coherence:ipc}(a)), and a fortiori $A\der^{\IPC}_{\mathrm{adm}}\bot$. \qed

% ---------------------------------------------------------------------
\subsection{Proof of \texorpdfstring{\cref{thm:coherence:ipc}}{the intuitionistic theorem}}\label{app:coherence:ipc}

\emph{(a)} Glivenko's theorem \citep{glivenko1929quelques} gives $\vdash_{\CPC}\neg\psi$ iff $\vdash_{\IPC}\neg\psi$. With $\psi=\bigwedge A$ and the deduction theorem (in both logics), $A\der_{\CPC}\bot$ iff $\vdash_{\CPC}\neg\bigwedge A$ iff $\vdash_{\IPC}\neg\bigwedge A$ iff $A\der_{\IPC}\bot$ (for $A=\emptyset$ neither side holds). An intermediate $L$ is sandwiched between $\IPC$ and $\CPC$.

\emph{(b)} $\CPC$ is a structural proper extension of $\IPC$, coherent on every classically consistent context; by (a) it has the same coherence data. So $\IPC$ is never a maximal coherent structural hypothesis.

\emph{(c)} We use the following facts about Jankov formulas \citep{jankov1968construction}; see \citet[ch.~9]{chagrov1997modal}. There is an infinite family $B_1,B_2,\dots$ of finite subdirectly irreducible Heyting algebras with $B_i\notin\mathbf{SH}(B_j)$ for $i\neq j$, and each $B_i$ has a formula $\chi(B_i)$ such that a Heyting algebra $B$ refutes $\chi(B_i)$ iff $B_i\in\mathbf{SH}(B)$. At most one $B_i$ is the two-element algebra $\mathbf2$ (two copies would violate the antichain condition); drop it, so $|B_i|>2$ for all $i$. Let $L_k$ be the least structural consequence relation containing $\der_{\IPC}$ and the axioms $\step\chi(B_i)$ for $i\le k$, and $L_\omega=\bigcup_kL_k$.
\begin{itemize}[leftmargin=*]
\item \emph{Strictly increasing.} The matrix consequence $\models_{B_{k+1}}$ of $(B_{k+1},\{1\})$ is structural and contains $\der_{\IPC}$, by soundness of $\IPC$ for Heyting algebras. $B_{k+1}$ validates $\chi(B_i)$ for $i\le k$, since $B_i\notin\mathbf{SH}(B_{k+1})$. So $L_k\subseteq{\models_{B_{k+1}}}$, whereas $B_{k+1}\in\mathbf{SH}(B_{k+1})$ refutes $\chi(B_{k+1})\in L_{k+1}(\emptyset)$.
\item \emph{Intermediate.} $\mathbf{SH}(\mathbf2)$ consists of $\mathbf2$ and the trivial algebra, and $B_i$ is neither, so $\mathbf2$ validates every $\chi(B_i)$. Hence every $L_k\subseteq{\models_{\mathbf2}}=\CPC$. By \cref{prop:coherence:chains}, $L_\omega$ is a consequence relation; it is structural as a union of structural relations, and $\IPC\subseteq L_\omega\subseteq\CPC$.
\item \emph{Impossibility.} By (a), for every $k\le\omega$ the set of $L_k$-coherent finite contexts is the set of classically consistent ones. So whether $\Ac$ is fixed, or designations are target-dependent data, the environment can present one and the same designation stream $\delta$ for every target in $\{L_k:k\le\omega\}$. Folding $\delta$ into the learner turns any learner for this class into a text learner for the chain $L_1\subsetneq L_2\subsetneq\cdots$ together with its union, which is impossible by the locking-sequence argument of \cref{cor:coherence:limitpoints}.
\end{itemize}

\emph{(d)} The target is $\langle G\rangle$ for a finite set $G$ of sequents. Since $D_n$ is contained in the structural target, so is $\langle D_n\rangle$. Once the text has shown every sequent of $G$, $\langle D_n\rangle\supseteq\langle G\rangle$. From then on every conjecture is correct (its syntactic index, e.g.\ $D_n$, may keep changing), which is BC-identification. \qed

% ---------------------------------------------------------------------
\subsection{Proof of \texorpdfstring{\cref{prop:coherence:complete}}{the complete-theories proposition}}\label{app:coherence:complete}

\emph{(a)} If a consistent deductively closed $T'\supsetneq T$ contains $\varphi\notin T$, then $\neg\varphi\in T\subseteq T'$ by completeness of $T$, so $T'$ is inconsistent. For learning, run the learner of \cref{thm:coherence:cpc}(b): for $j<i^*$, either $h_j\not\supseteq T$ and the text refutes it, or $h_j\supsetneq T$ is inconsistent and the datum $\emptyset$ rejects it at once.

\emph{(b)} Let $h$ be parameter-structural, $\exists$-closed, $h\supseteq{\der_T}$ and $\emptyset\nder_h\bot$. Suppose $\Gamma\der_h\varphi$ but $T\cup\Gamma\nvdash\varphi$, with parameters $\bar c=c_1,\dots,c_n$. Since the parameters do not occur in $T$, $T\nvdash\forall\bar x(\bigwedge\Gamma(\bar x)\to\varphi(\bar x))$, so by completeness $T\vdash\exists\bar x\,\theta(\bar x)$ for $\theta(\bar x):=\bigwedge\Gamma(\bar x)\wedge\neg\varphi(\bar x)$. Take fresh parameters $\bar d$. By parameter-structurality $\Gamma(\bar d)\der_h\varphi(\bar d)$. Since $h\supseteq{\der_T}$ is closed under cut and $\theta(\bar d)\der_T\gamma$ for each $\gamma\in\Gamma(\bar d)$, we get $\theta(\bar d)\der_h\varphi(\bar d)$; also $\theta(\bar d)\der_T\neg\varphi(\bar d)$ and $\varphi(\bar d),\neg\varphi(\bar d)\der_T\bot$, so $\theta(\bar d)\der_h\bot$. For $0\le i\le n$ let $\theta_i:=\exists x_{i+1}\cdots\exists x_n\,\theta(d_1,\dots,d_i,x_{i+1},\dots,x_n)$, so $\theta_n=\theta(\bar d)$ and $\theta_{i-1}=\exists x_i\,\psi(x_i)$ with $\psi(d_i)=\theta_i$. If $\theta_i\der_h\bot$, then $\exists$-elimination applied to $\theta_{i-1}\der_h\theta_{i-1}$ and $\theta_{i-1},\psi(d_i)\der_h\bot$ (monotonicity), with $d_i$ not in $\theta_{i-1}$ or $\bot$, gives $\theta_{i-1}\der_h\bot$. By induction $\theta_0\der_h\bot$, and since $\emptyset\der_T\theta_0$, cut gives $\emptyset\der_h\bot$, a contradiction. (With no parameters, $T\vdash\neg(\bigwedge\Gamma\to\varphi)$ and the last step is immediate.) So $h\subseteq{\der_T}$, i.e.\ $h={\der_T}$. Learning is as in (a).

\emph{(c)} Terms are integer polynomials in the parameters. For a set $X$ of sentences, let $\mathrm{cl}_i(X)$ be the least set containing $X$, closed under $\Cn_{\RCF}$ and under the instances of rule $i$ ($t\cdot t=1+1\step t>0$ for $i=1$; $t^3-4t+1=0\step t>1$ for $i=2$). Substituting terms for parameters maps $\Cn_{\RCF}$-inferences to $\Cn_{\RCF}$-inferences and rule instances to rule instances, so $\mathrm{cl}_i$ is parameter-structural, and it is the least parameter-structural consequence relation $h_i$ containing $\der_{\RCF}$ and the rule. Because the parameters do not occur in $\RCF$ and $\RCF$ is complete with model $\R$, $\RCF\cup Y\vdash\psi$ for parameter sentences iff $\R$ satisfies the universal closure of $\bigwedge Y_0\to\psi$ for some finite $Y_0\subseteq Y$.

\emph{$h_1$.} An instance fires from $\Cn_{\RCF}(\emptyset)$ only if $t(\bar x)^2=2$ identically on $\R$, which no integer polynomial satisfies (a constant integer $n$ has $n^2\neq2$, and a non-constant $t$ has non-constant $t^2$). So $h_1(\emptyset)=\Cn_{\RCF}(\emptyset)$, which is consistent. The instance with $t=c$ gives $c\cdot c=1+1\der_{h_1}c>0$, which $\RCF$ does not prove ($c=-\sqrt2$). In the context $\{c\cdot c=1+1\}$, $\RCF$ proves $(-c)\cdot(-c)=1+1$, the rule gives $-c>0$ and $c>0$, and $\RCF$ derives $\bot$.

\emph{$h_2$.} The cubic $f(x)=x^3-4x+1$ has no rational root ($f(\pm1)\neq0$), so it is irreducible over $\mathbb Q$. Its discriminant is $-4(-4)^3-27=229$, not a square, so its Galois group is $S_3$. Its real roots are $r_1\approx-2.115$, $r_2\approx0.254$, $r_3\approx1.861$. For a root $r$, $\mathbb Q(r)$ has degree $3$ and is not normal (else the Galois group would have order $3$), so it contains no other root (if it contained two roots it would contain the third, since the roots sum to $0$, and be the splitting field). Now:
\begin{itemize}[leftmargin=*]
\item From $\emptyset$, an instance fires only if $f(t(\bar x))=0$ for all real $\bar x$. Then $t$ is a continuous map from the connected $\R^m$ into the finite set of roots, hence constant, hence an integer root of $f$, which does not exist. So $h_2(\emptyset)=\Cn_{\RCF}(\emptyset)$.
\item Let $X=\{f(c)=0\}$ and $Y=\Cn_{\RCF}(X\cup\{c>1\})$. The instance with $t=c$ shows $Y\subseteq h_2(X)$. Conversely, $Y$ is closed under the rule: if $f(t(c,\bar d))=0\in Y$, then $f(t(r_3,\bar d))=0$ for all real $\bar d$ (the only root $>1$ is $r_3$); by continuity and connectedness $t(r_3,\bar d)$ is constant in $\bar d$, equal to $t(r_3,\bar0)\in\mathbb Z[r_3]\subseteq\mathbb Q(r_3)$, hence equal to $r_3$; so $\R$ satisfies $\forall c\,\forall\bar d\,(f(c)=0\wedge c>1\to t(c,\bar d)>1)$ and $t>1\in Y$. Hence $h_2(X)=Y$, which is consistent, while $X\der_{h_2}c>1$ is not $\RCF$-valid ($c=r_2$). So $h_2$ is unsound, and coherent on both $\emptyset$ and $X$.
\item In the context $\{f(c)=0,\ c<0\}$ the rule gives $c>1$, and $\RCF$ derives $\bot$. \qed
\end{itemize}

\emph{Repair note.} The first version claimed, from (a), that complete theories are pinned ``as cleanly as $\CPC$'' in the user's rule-level setting. The verifier's counterexamples (c) refuted that; (b) is the repair. The algebraic facts about $f$ were also checked with a computer algebra system.

% ---------------------------------------------------------------------
\subsection{Proof of \texorpdfstring{\cref{thm:coherence:popper}}{the Popperian theorem}}\label{app:coherence:popper}

\emph{(a)} If $\pi$ is true, $\PA+\pi\subseteq\Th(\N)$ is consistent. If $\pi$ is false, $\neg\pi$ is a true $\Sigma_1$ sentence, so $\PA\vdash\neg\pi$ by $\Sigma_1$-completeness, and $\PA+\pi$ is inconsistent.

\emph{(b)} A true $\sigma$ gives $\PA+\sigma\subseteq\Th(\N)$. By G\"odel's second theorem $\PA+\neg\Con(\PA)$ is consistent, although $\neg\Con(\PA)$ is a false $\Sigma_1$ sentence.

\emph{(c)} On a $\Pi_1$ sentence $\pi$ the learner says ``true'' at stage $t$ unless a $\PA$-proof of $\neg\pi$ with code $\le t$ exists. By (a) and soundness of $\PA$, such a proof exists iff $\pi$ is false; so the verdict changes at most once and converges to the truth. On a $\Sigma_1$ sentence $\sigma$ it says ``false'' unless a $\PA$-refutation of $\neg\sigma$ (equivalently a proof of $\sigma$) with code $\le t$ exists; apply (a) to the $\Pi_1$ sentence $\neg\sigma$.

\emph{(d)} At stage $t$ the learner accepts $\varphi_n$ iff no $\PA$-proof with code $\le t$ refutes the conjunction of $\varphi_n$ with the sentences among $\varphi_0,\dots,\varphi_{n-1}$ accepted at stage $t$. By induction on $n$: once the statuses of $\varphi_0,\dots,\varphi_{n-1}$ have stabilized, the status of $\varphi_n$ is a $\Pi_1$ question (consistency of a fixed finite extension of $\PA$), answered by proof search with at most one change. So all statuses converge. In the limit $\varphi_n$ is accepted iff it is consistent with $\PA$ plus the accepted earlier sentences: this is Lindenbaum's construction along the enumeration, so the limit is a complete consistent extension of $\PA$, and it depends on the enumeration. If $\varphi_0=\neg\Con(\PA)$, no refutation of $\PA+\neg\Con(\PA)$ exists (G\"odel~II), so $\varphi_0$ is accepted forever and the limit contains a false $\Sigma_1$ sentence. On the enumeration $\Con(\PA)\wedge0{=}0,\ \neg\Con(\PA),\ \Con(\PA),\dots$, the first sentence is consistent with $\PA$ (it is true) and stays accepted; the second is inconsistent with it, which proof search finds, so it is rejected.

\emph{(e)} On a computable presentation, the learner's guesses about each sentence form a computable function $g(\varphi,n)$. If they converged to the truth value on every $\Sigma_2$ sentence, then by Shoenfield's limit lemma \citep{shoenfield1959degrees} the set of true $\Sigma_2$ sentences would be $\Delta_2$; but it is $\Sigma_2$-complete, hence not $\Delta_2$, by Post's hierarchy theorem. \qed

\emph{Repair note.} The first version of (d) said that the naive learner accepts whichever of $\Con(\PA)$ and $\neg\Con(\PA)$ comes first; the verifier's enumeration above refutes that, and (d) now states the Lindenbaum form.

% ---------------------------------------------------------------------
\subsection{Proof of \texorpdfstring{\cref{prop:coherence:costs}}{the weak-coherence cost proposition}}\label{app:coherence:costs}

\emph{(a)} If limits exist on $C$, then for $\varphi\in C$: $\lim_tP_t(\varphi)=1$ iff $\forall q\in\mathbb Q^+\,\forall s\,\exists t\ge s\ P_t(\varphi)>1-q$, a $\Pi_2$ condition. The set of true $\Sigma_2$ sentences is $\Sigma_2$-complete, so it is not $\Pi_2$ (else it would be $\Delta_2$), and no convergent computable credences have it as their success set.

\emph{(b)} By the convention on empty quantifier blocks every $\Pi_1$ sentence is $\Pi_2$. If weakly coherent credences gradually verified all $\Pi_2$ sentences, they would be $\Pi_1$-convergent, and \cref{thm:coherence:trilemma} would give a true $\Pi_2$ sentence with $\liminf_tP_t=0$, contradicting $P_t\to1$ on true $\Pi_2$ sentences.

\emph{(c)} Enumerate the sentences of $\Sigma_2$ form, $\sigma=\exists x\forall y\,R(x,y)$ with $R$ in $\Delta_0$ (quantifier blocks coded as single variables; an empty block is a dummy variable), with indices $\mathrm{ind}(\sigma)$. At stage $t$:
\begin{itemize}[leftmargin=*]
\item $x_t(\sigma)$ is the least $x\le t$ with $R(x,y)$ for all $y\le t$, or $t+1$ if there is none;
\item $c_t(\sigma)$ is the last stage $s\le t$ with $x_s(\sigma)\neq x_{s-1}(\sigma)$, or $-1$ if there is none;
\item $a_t(\sigma):=t-\max(\mathrm{ind}(\sigma),c_t(\sigma))$, the age of the current witness, for $t\ge\mathrm{ind}(\sigma)$;
\item $I_t$ is the symmetric set of pairs $(\sigma,\sigma')$ of $\Sigma_2$ sentences with indices $\le t$ that have a $\PA$-proof of $\neg(\sigma\wedge\sigma')$ with code $\le t$ (including $\sigma'=\sigma$);
\item $P_t(\sigma):=\min\bigl(1-2^{-a_t(\sigma)},\ \min\{2^{-a_t(\sigma')}:(\sigma,\sigma')\in I_t,\ a_t(\sigma')\ge a_t(\sigma)\}\bigr)$.
\end{itemize}
Set $P_t:=0$ on sentences not of $\Sigma_2$ form, and on $\sigma$ before stage $\mathrm{ind}(\sigma)$.

\emph{Weak coherence.} If $\PA\vdash\psi\to\neg\varphi$ and one of them is not of $\Sigma_2$ form, it has credence $0$ and the sum is at most $1$. Otherwise $(\psi,\varphi)\in I_t$ from some stage on. At such a stage, say $a_t(\psi)\le a_t(\varphi)$; then $P_t(\psi)\le2^{-a_t(\varphi)}$ and $P_t(\varphi)\le1-2^{-a_t(\varphi)}$, so the sum is at most $1$.
\emph{False $\sigma$.} Every $x$ is eventually refuted, so $x_t(\sigma)\to\infty$ and changes infinitely often. At each change $a_t(\sigma)=0$ and $P_t(\sigma)\le1-2^0=0$. So $P_t(\sigma)\not\to1$.
\emph{True $\sigma$.} Let $x^*$ be its least witness. Once every $x<x^*$ is refuted and $t\ge x^*$, $x_t(\sigma)=x^*$ forever. Let $T_0=\max(\mathrm{ind}(\sigma),\text{last change})$; for $t>T_0$, $a_t(\sigma)=t-T_0\to\infty$. A competitor $\sigma'$ with $(\sigma,\sigma')\in I_t$ is $\PA$-incompatible with the true $\sigma$, hence false since $\PA$ is sound. If $\mathrm{ind}(\sigma')>T_0$ then $a_t(\sigma')\le t-\mathrm{ind}(\sigma')<a_t(\sigma)$. The finitely many competitors with index $\le T_0$ are false, so each changes its witness at some stage $c>T_0$, after which $a_t(\sigma')\le t-c<a_t(\sigma)$. So from some stage on no competitor enters the inner minimum, and $P_t(\sigma)=1-2^{-a_t(\sigma)}\to1$.
\emph{Non-convergence.} By (a), these credences cannot converge on every $\Sigma_2$ sentence. \qed

The construction is due to L3's referee and was sanity-checked by script (\nolinkurl{L3-scripts/sigma2_weak_coherence.py}). It assigns $0$ to every sentence not of $\Sigma_2$ form, so it is dogmatic on the strictly $\Pi_2$ truths, consistent with \cref{thm:coherence:trilemma}.

% ---------------------------------------------------------------------
\subsection{Proofs for \texorpdfstring{\cref{sec:coherence:carnap}}{Section 6.4}: Carnap's problem}\label{app:coherence:carnap}

For finite $\Gamma,\Delta$ write $U(\Gamma,\Delta)=\{v:v[\Gamma]=1,\ v[\Delta]=0\}$; these sets form a clopen basis of $2^{\Fm}$, and a valuation satisfies the sequent $\Gamma\step\Delta$ iff it lies outside $U(\Gamma,\Delta)$. So $\Val(S)$ is closed for every set $S$ of sequents.

\paragraph{\Cref{lem:coherence:duality}.} \emph{(a)} If $v\notin\overline V$, some basic neighbourhood $U(\Gamma,\Delta)\ni v$ misses $V$; then $\Gamma\models^{\mathrm m}_V\Delta$ and $v$ violates it. If $v\in\overline V$ violates a valid $\Gamma\step\Delta$, then $U(\Gamma,\Delta)$ is a neighbourhood of $v$, so it contains some $u\in V$, which violates the sequent too; contradiction.

\emph{(b)} \emph{Finitary:} if $\varphi\in C_V(X)$, the closed sets $V\cap\{v:v(\varphi)=0\}$ and $\{v:v(x)=1\}$, $x\in X$, have empty intersection; by compactness a finite subfamily already does, so $\varphi\in C_V(X_0)$ for a finite $X_0\subseteq X$. Hence a valuation respects every valid finite-premise single-conclusion sequent iff $C_V(v)=v$ (identifying $v$ with its true-set). \emph{$\mathrm{Fix}(C_V)=V^\cap$:} a fixed point $X$ equals $\bigcap\{v\in V:X\subseteq v\}$, an intersection of members of $V$ (the empty intersection being $\Fm$); conversely, for $X=\bigcap W$ with $W\subseteq V$, $C_V(X)\subseteq\bigcap W=X$.

\emph{(c)} An intersection of classical theories is a theory, and $\Fm$ is the empty intersection. Conversely, a consistent theory $T$ is the intersection of the Boolean valuations extending it: if $\varphi\notin T$, completeness gives a Boolean $v\supseteq T$ with $v(\varphi)=0$. \qed

\paragraph{\Cref{thm:coherence:telltale}.} \emph{(a)} Each schema enforces one row condition. For $\neg$: $\varphi,\neg\varphi\step$ excludes both true and $\step\varphi,\neg\varphi$ excludes both false, so $v(\neg\varphi)=1-v(\varphi)$. For $\wedge$: the two eliminations give $v(\varphi\wedge\psi)\le\min$, introduction gives $\ge$. For $\vee$: the introductions give $v(\varphi\vee\psi)\ge\max$, and $\varphi\vee\psi\step\varphi,\psi$ gives $\le$. For $\to$: $\step\varphi,\varphi\to\psi$ gives $v(\varphi\to\psi)=1$ when $v(\varphi)=0$; $\psi\step\varphi\to\psi$ gives it when $v(\psi)=1$; modus ponens gives $v(\varphi\to\psi)=0$ when $v(\varphi)=1,v(\psi)=0$. So $v$ satisfies all instances iff it is a homomorphism into the two-element Boolean algebra. (With $\bot,\top$ in the language, $\bot\step$ and $\step\top$ fix their values.) $\BV=\Val(\TTs)$ is closed, and $\models^{\mathrm m}_{\BV}$ determines $\BV$ by \cref{lem:coherence:duality}(a).

\emph{(b)} For any substitution $\sigma$ and sequent $s=\Gamma\step\Delta$, $v$ satisfies $\sigma s$ iff $v\circ\sigma$ satisfies $s$. Let $v\in V$. Every instance of a schema in $\TTs$ is $\sigma(\TTs_j(p,q))$ for some $\sigma$ and $j$, and $v\circ\sigma\in V$ satisfies $\TTs_j(p,q)$; so $v$ satisfies all instances and $V\subseteq\Val(\TTs)=\BV$. For equality, take any $v\in V$ (there is one, as $V\neq\emptyset$) and any $u\in\BV$. Let $\sigma(p)=p_0\vee\neg p_0$ if $u(p)=1$ and $\sigma(p)=\neg(p_0\vee\neg p_0)$ otherwise. Then $v\circ\sigma\in V\subseteq\BV$ is Boolean and agrees with $u$ on atoms ($v$ is Boolean, so $v(p_0\vee\neg p_0)=1$), hence $v\circ\sigma=u$ and $u\in V$.

\emph{(c)} The learner that outputs $\BV$ once the eleven sequents $\TTs(p,q)$ and the coherence datum have appeared is correct by (b), and it never needs to change its mind. $V=\emptyset$ satisfies every sequent vacuously, so without the coherence datum no finite set of positive data excludes it.

\emph{(d)} \emph{Pointwise refutation.} A text for $L_{\BV}:={\models^{\mathrm m}_{\BV}}$ eventually presents every instance of $\TTs$, and by (a) every non-Boolean valuation violates one. \emph{No tell-tale.} Fix a Boolean $u_0$, and for a formula $\chi$ let $v'_\chi$ agree with $u_0$ except at $\chi$. Then $v'_\chi(\neg\chi)=u_0(\neg\chi)=1-u_0(\chi)=v'_\chi(\chi)$, so $v'_\chi\notin\BV=\overline{\BV}$. $\BV\cup\{v'_\chi\}$ is closed, and by \cref{lem:coherence:duality}(a) its relation $L_\chi$ satisfies $\Val(L_\chi)=\BV\cup\{v'_\chi\}\neq\BV$, so $L_\chi\subsetneq L_{\BV}$. Given a finite $T\subseteq L_{\BV}$, choose $\chi$ occurring in no sequent of $T$; then $v'_\chi$ agrees with $u_0$ on every formula $T$ mentions and satisfies $T$, so $T\subseteq L_\chi$. \emph{Non-learnability.} Suppose a learner, computable or not, BC-identifies (a fortiori if it EX-identifies) every meaning in a class containing $\BV$ and all $\BV\cup\{v'_\chi\}$ from texts of their relations, plus the coherence datum $\pos{\ }{\ }$, which is in bounds for every member and so carries no information. By the locking-sequence argument \citep{blum1975toward,angluin1980inductive}, which needs no effectivity, there is a finite $\sigma$ of $L_{\BV}$-data such that the learner conjectures $\BV$ on every extension of $\sigma$ by $L_{\BV}$-data. Choose $\chi$ not mentioned in $\sigma$; then $\sigma$ consists of $L_\chi$-data, and on a text for $L_\chi$ beginning with $\sigma$ the learner conjectures $\BV\neq\BV\cup\{v'_\chi\}$ forever. \emph{Density.} A Boolean valuation is determined by its values on atoms, and the map from atom assignments to $\BV$ is a homeomorphism from $2^\omega$, which has no isolated points. So $\BV\setminus\{u_0\}$ is dense in $\BV$; by \cref{lem:coherence:duality}(a) it has the same multiple-conclusion relation, and a basic open set meets $\BV$ iff it meets $\BV\setminus\{u_0\}$, so the two meanings have the same in-bounds positions. Composing any Boolean $v$ with the constant substitution along $u_0$ yields $u_0$, so $\BV\setminus\{u_0\}$ is not structural. \qed

\emph{Repair note.} The first version of (d) ended ``identification in the limit still holds via the text''. That was false and has been retracted; (d) now proves non-identifiability. The Lean formalization proves (d) for BC-identification from texts with pauses and arbitrary learners.

\paragraph{\Cref{prop:coherence:compositional}.} Row by row, with designated set $\{1\}$. \emph{$\neg$:} explosion $p,\neg p\step q$ at $p=1$, $q=0$ forces $f_\neg(1)=0$; then $p\step\neg\neg p$ at $p=1$ forces $f_\neg(0)=1$. \emph{$\wedge$:} $p\wedge q\step p$ forces $f_\wedge(0,b)=0$, $p\wedge q\step q$ forces $f_\wedge(a,0)=0$, and $p,q\step p\wedge q$ forces $f_\wedge(1,1)=1$. \emph{$\vee$:} the introductions force $f_\vee(a,1)=f_\vee(1,b)=1$, and disjunctive syllogism at $p=q=0$ (where $\neg p=1$) forces $f_\vee(0,0)=0$. \emph{$\to$:} modus ponens forces $f_\to(1,0)=0$; $q\step p\to q$ forces $f_\to(a,1)=1$; $\neg p\step p\to q$ at $p=0$ forces $f_\to(0,b)=1$. So only the standard matrix can validate all twelve rules, and it does. A brute-force check over all $4\cdot16^3$ matrices confirms that exactly one matrix validates them (\nolinkurl{T2-checks/matrix_check.py}). \qed

% ---------------------------------------------------------------------
\subsection{Proofs for \texorpdfstring{\cref{sec:coherence:tonk}}{Section 6.5}: tonk and conservativity}\label{app:coherence:tonk}

\paragraph{\Cref{prop:coherence:conservativity}.} Intuitionistic implication with $\neg$-introduction, $\neg$-elimination and $\neg\neg$-elimination is a standard natural-deduction system for the $\{\to,\neg\}$ fragment of classical logic \citep{prawitz1965natural}, so it is coherent and proves Peirce's law $((p\to q)\to p)\to p$. Peirce's law is not intuitionistically valid: in the two-world Kripke model $w\le w'$ with $p$ true only at $w'$ and $q$ true nowhere, $p\to q$ fails at both worlds, so $(p\to q)\to p$ holds at $w$ vacuously, while $p$ fails at $w$. So the extension is not conservative over $\IPC_\to$.
\emph{$\neg\neg$-elimination alone is conservative.} Let $\tau$ erase negations ($\tau(\neg\varphi)=\tau(\varphi)$, $\tau$ commuting with $\to$, identity on atoms). It maps each instance of $\neg\neg\varphi\step\varphi$ to the identity step $\tau\varphi\step\tau\varphi$ and each instance of an $\IPC_\to$ rule to an instance of the same rule, and it fixes $\neg$-free formulas. So a derivation of a $\neg$-free sequent in the extension maps to an $\IPC_\to$-derivation of the same sequent.
\emph{Conservativity and coherence.} If the base is non-trivial and coherence means non-triviality on the old vocabulary, a conservative extension proves no new old-vocabulary sequent, so it is non-trivial; likewise if $\bot$ with ex falso is in the base. If $\bot$ is new and has no rules, $\IPC_\to+\{\step\bot\}$ adds no $\bot$-free theorem and so is conservative, yet derives $\bot$. \qed

\paragraph{\Cref{thm:coherence:complexity}.} \emph{(a)} ``$R\nder\bot$'' says that no finite derivation in $R$ ends in $\bot$, a $\Pi_1$ condition. For hardness, given $e$ build a finite system $R_e$ that derives the successive configurations of the computation of $\varphi_e(e)$ (elementary formal systems simulate Turing machines; \citealt{post1943formal,smullyan1961theory}), and add the rule ``halting configuration $\step\bot$''. Then $R_e\nder\bot$ iff $\varphi_e(e)$ diverges, which reduces the $\Pi_1$-complete complement of the halting set.

\emph{(b)} \emph{Membership:} ``for every old-vocabulary $s$ and every derivation of $s$ in $B+N$ there is a derivation of $s$ in $B$'' is $\Pi_2$. \emph{Hardness}, from the $\Pi_2$-complete set $\mathrm{Tot}=\{e:\varphi_e\text{ total}\}$. Let $B$ be a universal elementary formal system that derives $T(\bar x,\bar e)$ iff $\varphi_e(x)\!\downarrow$ (unary numerals $\bar x,\bar e$) and derives $\mathrm{Num}(\bar n)$ for exactly the unary numerals, designed so that $T$-strings never occur as premises of its rules. Its productions must also match only predicate-tagged strings over the old alphabet; in pure Post-string format a projection such as $xy\step y$ could otherwise strip a new symbol, while in an elementary formal system with predicates this is automatic. Let $N_e$ add a new symbol $c$ with the rules $\step c\,\bar e$ and $c\,y,\ \mathrm{Num}(x)\step T(x,y)$ (schematic in $x,y$). The only derivable $c$-string is $c\,\bar e$, so the new old-vocabulary theorems are exactly $T(\bar n,\bar e)$ for $n\in\N$, and since $T$-strings are never premises nothing further follows. So the old-vocabulary theorems of $B+N_e$ are $\mathrm{Thm}(B)\cup\{T(\bar n,\bar e):n\in\N\}$, and $B+N_e$ is conservative over $B$ iff $\varphi_e(n)\!\downarrow$ for every $n$, iff $e\in\mathrm{Tot}$.

\emph{(c)} A set is decidable in the limit by a computable learner iff it is $\Delta_2$ \citep{shoenfield1959degrees,putnam1965trial,gold1965limiting}; a $\Pi_1$ set is decided in the limit with at most one mind change; a $\Pi_2$-complete set is not $\Delta_2$. Refutability in the limit corresponds to $\Pi_2$ \citep{kelly1996logic}.

\emph{(d)} If $B$ is decidable, ``every old-vocabulary theorem of $B+N$ is a theorem of $B$'' quantifies universally over derivations and then checks a decidable property: $\Pi_1$. \qed

% ---------------------------------------------------------------------
\subsection{Proof of \texorpdfstring{\cref{thm:coherence:residue}}{the residue theorem}}\label{app:coherence:residue}

The characterization of $\Alt$ as the never-refuted $\mathcal G$-closed hypotheses is \cref{thm:coherence:silent}(ii), restricted to $\mathcal G$-closed $h$.

\emph{(i)} A structural $h\supseteq\Cn_2$ is $\Cn_2$ or $\CFm$ by \cref{thm:coherence:post}, and $\CFm$ is incoherent on every context, so on the non-empty $\Ac$.

\emph{(ii)} \emph{$\supseteq$:} if $h$ is structural with $\IPC\subseteq h\subseteq\CPC$ and $A$ is classically consistent, then $A\nder_{\CPC}\bot$, so $A\nder_h\bot$. \emph{$\subseteq$:} let $h\supseteq\IPC$ be structural and suppose $\Gamma\der_h\varphi$ with $\Gamma$ finite and $\Gamma\not\models\varphi$. Take a Boolean $v$ with $v[\Gamma]=1$, $v(\varphi)=0$, and let $\sigma_v$ map each atom to $\bot\to\bot$ or $\bot$ according to $v$. As shown in the proof of \cref{prop:coherence:adm}(c), $\IPC$ proves every classically true variable-free formula and refutes every classically false one. So $\IPC$ proves each member of $\sigma_v\Gamma$ and proves $\neg\sigma_v\varphi$. By structurality $\sigma_v\Gamma\der_h\sigma_v\varphi$; by cut with the $\IPC$-theorems $\sigma_v\Gamma$, $\emptyset\der_h\sigma_v\varphi$; with $\emptyset\der_h\neg\sigma_v\varphi$ and $\sigma_v\varphi,\neg\sigma_v\varphi\der_{\IPC}\bot$, $\emptyset\der_h\bot$. By monotonicity $h$ is incoherent on every $A\in\Ac$, which is non-empty. So every coherent structural $h\supseteq\IPC$ lies inside $\CPC$.

\emph{(iii)} With $\mathcal G$ deductive closure and $\Ac=\{\emptyset\}$, the hypotheses above the target are the deductively closed theories containing $\PA$, and $\emptyset$-coherence is consistency. Designating $\{\Con(\PA)\}$ removes $\PA+\neg\Con(\PA)$. If every finite set of true sentences is designated, a theory containing a false $\varphi$ is incoherent on $\{\neg\varphi\}$, and a theory $T\subseteq\Th(\N)$ is coherent on every true context; so exactly the unsound ones are removed. The theories $T_k$ of \cref{thm:coherence:turing} contain $\PA$ and are sound, so they remain, and by that theorem text plus true-context coherence cannot discriminate among them. \qed

\emph{Repair note.} Part (ii) was strengthened from ``$\supseteq$'' to equality after verification, with the $\sigma_v$ argument above; part (iii) was made precise ($\mathcal G$ and $\Ac$ specified).
